\documentclass[a4paper, 12pt]{article}
\usepackage[UTF8]{ctex}
\usepackage{hyperref}

\begin{document}
    \title{事已至此，还是回家睡觉吧,吃点好的}
    \author{}
    \date{\today , \; 记}
    \maketitle

    \section{我的概况}
    我目前大三了，大概也是被22年高考数学新国二卷狙击了。
    英语144，语文124，但是数学炸裂，十分可惜，就被送来桂林的山里修仙了。
    但专业是计算机相关的，挺满足了。
    
    本科的前两年都在打传统算法竞赛，最近半年也在搞AI算法，在CS里算是和数学强相关的了。
    不过我也只能给给，从计算机角度，对数学专业的一些理解。

    \section{我对数学的一些理解}
    我对中学阶段的数学理解，大概印象是抱着一些初等函数做灵能符号变换，抱着块烂泥胚子雕花吧，没什么好感。
    这个观点和《上海交通大学生存手册》里写的，不谋而合吧。
    
    本科专业细分之后，每个学科有自己的数学工具，事情才开始变得有意思。
    像计算机，线代、数论、图论、多项式生成函数和组合数学应该是传统CS的风格代表。
    当然近十年AI兴起后，概率论也算是代表了吧。微积分、偏微分方程、复变函数之类的，CS反而用的不多。
    oi-wiki这个网站的数学板块挺全的，没事可以去看看。

    \begin{samepage}
        \noindent CS里对数学的运用，我觉得能抽象为以下三个状态：
        \begin{enumerate}
            \item 对我的研究对象，做一个恰当的数学表达
            \item 在数学表达的基础上（特别是代数表达）推演数学结构和性质
            \item 回到研究对象，使用适当的数学工具对解析出的数学结构性质处理
        \end{enumerate}
    \end{samepage}

    这个观点取自南京大学的蒋炎岩教授对数学在信奥竞赛中的理解。
    只不过我觉得，推广到CS的多个领域，也同样适用。

    举个例子，AI算法中的模型训练，按以上的理论，抽象成以下三个状态：
    \begin{samepage}
        \begin{enumerate}
            \item 对我的模型性能，采用一个实变函数进行刻画，将其称为损失函数。函数值越小，模型性能越好。
            \item 将模型性能转化为一个函数后，推演出该函数可微分性、凹凸性、极值等函数性质。
            \item 利用凸优化、数值优化体系里的数学工具，处理这些性质，找到令损失函数值尽可能小的模型参数。
        \end{enumerate}
    \end{samepage}

    一般而言，CS会更注重在3上。数学专业会更注重1和2。
    其实都是做数学吧，就看你喜欢哪种风格多一点了。

    \section{我对研究生的一些理解}
    \href{https://jyywiki.cn/Letter.md}{这是蒋炎岩老师写的一封对本科生的“科研劝退信”，决定读研前可以看看。}
    说说我的补充理解吧。草台班子氛围盛行的环境，遇上一个人品过得去的导师，已是
    人生未来三年一大幸事，老板的学术能力，我已不敢奢求了。
    要是能遇到手把手教你配环境写代码的导，那大抵就是再生父母了吧。
    
    说说深度学习方向的科研风格吧。一般的科研流水线兵团，就是导师划定领域，博士输出idea，硕士写代码验证。
    dl领域的idea挺廉价的，主要依赖代码实现去验证，要是不会coding，大概率就变成实验室的负资产了。
    这主要还是个实验大于理论的领域吧，但也有很多人尝试去完善理论，特别是做模型可解释性方向，很数学，也很有挑战性。
    
    哦哦，听闻华东师范的数学专业，挺猛的，男女比例也挺好的。学校在上海，是个好地方，纸醉金迷的，总比在山里好。
    
    
\end{document}